\documentclass[18pt]{beamer}

\usetheme{metropolis}
\usefonttheme{professionalfonts}


%% "Patch" to keep the text widths' of Metropolis similar to that of Madrid.
%\setbeamersize{text margin left=15pt, text margin right=15pt}
%\makeatletter
%\setbeamertemplate{title page}{
%\centering
%  \begin{minipage}[b][\paperheight]{.9\textwidth}
%    \ifx\inserttitlegraphic\@empty\else\usebeamertemplate*{title graphic}\fi
%    \vfill%
%    \ifx\inserttitle\@empty\else\usebeamertemplate*{title}\fi
%    \ifx\insertsubtitle\@empty\else\usebeamertemplate*{subtitle}\fi
%    \usebeamertemplate*{title separator}
%    \ifx\beamer@shortauthor\@empty\else\usebeamertemplate*{author}\fi
%    \ifx\insertdate\@empty\else\usebeamertemplate*{date}\fi
%    \ifx\insertinstitute\@empty\else\usebeamertemplate*{institute}\fi
%    \vfill
%    \vspace*{1mm}
%  \end{minipage}
%}
%\makeatother

% Modify the text widths' of section title pages. See `beamerinnerthememetropolis.dtx`.
% `\paperheight` requires adjustment after modifying the frame content placement. Not sure why.
\makeatletter
\setbeamertemplate{section page}{
  \centering
  \begin{minipage}[c][.95\paperheight]{.9\textwidth}
    \raggedright
    \usebeamercolor[fg]{section title}
    \usebeamerfont{section title}
    \insertsectionhead\\[-1ex]
    \usebeamertemplate*{progress bar in section page}
    \par
    \ifx\insertsubsectionhead\@empty\else%
      \usebeamercolor[fg]{subsection title}%
      \usebeamerfont{subsection title}%
      \insertsubsectionhead
    \fi
  \end{minipage}
  \par
  \vspace{\baselineskip}
}
\makeatother

% Make frametitles larger and place closer to center. See `beamerinnerthememetropolis.dtx`.
\setbeamerfont{frametitle}{size=\Large}
\makeatletter
% Theme default adds paddings of 2.2ex to all the four sides of frame titles
\newlength{\metropolis@frametitle@toppadding}
\newlength{\metropolis@frametitle@bottompadding}
\setlength{\metropolis@frametitle@toppadding}{3.5ex} 
\setlength{\metropolis@frametitle@bottompadding}{0ex} 
\setlength{\metropolis@frametitle@padding}{4ex} % Horizontal padding to left
\renewcommand{\metropolis@frametitlestrut@start}{
  \rule{0pt}{\metropolis@frametitle@toppadding +%
    \totalheightof{%
      \ifcsdef{metropolis@frametitleformat}{\metropolis@frametitleformat X}{X}%
    }%
  }%
}
\newcommand*\getlength[1]{\number#1}
\renewcommand{\metropolis@frametitlestrut@end}{%
  \ifnum\getlength{\metropolis@frametitle@bottompadding}>0%
    \rule[-\metropolis@frametitle@bottompadding]{0pt}{\metropolis@frametitle@bottompadding}%
  \fi%
}
\makeatother

% Move up the frame content a bit. See:
% https://github.com/matze/mtheme/blob/master/source/beamerinnerthememetropolis.dtx#L509-L523
% https://tex.stackexchange.com/questions/247826/beamer-full-vertical-centering
\makeatletter
\define@key{beamerframe}{c}[true]{% centered
  \beamer@frametopskip=0pt plus 0.5fill\relax% Default is `1fill`
  \beamer@framebottomskip=0pt plus 1.5fill\relax%
  \beamer@frametopskipautobreak=0pt plus .4\paperheight\relax%
  \beamer@framebottomskipautobreak=0pt plus .6\paperheight\relax%
  \def\beamer@initfirstlineunskip{}%
}
\makeatother


%% Packages
\usepackage[english]{babel}
\usepackage{graphicx}
\usepackage{mathtools} % For \coloneqq (among others)
\usepackage{subcaption}
\usepackage{bm}
\usepackage{accents}
\usepackage{pgfplots}
%\usepackage[normalem]{ulem}
%	\renewcommand{\ULthickness}{.6pt} % default is 0.4pt
%	\newcommand{\coloredUline}{\bgroup\markoverwith{\textcolor{lava}{\rule[-0.5ex]{2pt}{0.6pt}}}\ULon}
\usepackage{natbib}
	\bibliographystyle{abbrvnat}
	\setcitestyle{authoryear, open={(}, close={)}}
\usepackage{usebib}
\usepackage{bibentry} % For \nobibliography
%\nobibliography{references} % To cite inline without creating a bibliography
\newbibfield{journal}
\bibinput{references}


%% Beamer options.
\setbeamercovered{dynamic}
\setbeamercovered{invisible}
\beamertemplatenavigationsymbolsempty
%\setbeamertemplate{caption}{unnumbered}
\setbeamertemplate{footline}{}
%\setbeameroption{show notes on second screen}

%% Customize styles inside itemize environments
%\settowidth{\leftmargini}{\usebeamertemplate{itemize item}}
%\addtolength{\leftmargini}{-1.2\labelsep}
\setbeamerfont{itemize/enumerate subbody}{size=\normalsize} %to set the body size
\setbeamertemplate{itemize subitem}{\normalsize\raise1.25pt\hbox{\donotcoloroutermaths$\blacktriangleright$}}  % to set the symbol size

%% Custom environments
\newcommand{\defineTightItemizeSpacing}{%
	\setlength{\abovedisplayskip}{.25\baselineskip}%
	\setlength{\belowdisplayskip}{.25\baselineskip}%
}
\newenvironment{tightEquation}{%
	\defineTightItemizeSpacing%
	\begin{equation}
}{
	\end{equation} \ignorespacesafterend
}
\makeatletter
\newenvironment{tightEquation*}{%
	\defineTightItemizeSpacing%
	\begin{equation*}
}{
	\end{equation*} \ignorespacesafterend
}
\newcommand{\defineWithinItemizeSpacing}{%
	\setlength{\abovedisplayskip}{.3\baselineskip}%
	\setlength{\belowdisplayskip}{.3\baselineskip}%
}
\newenvironment{itemizedEquation}{%
	\defineWithinItemizeSpacing%
	\begin{equation}
}{
	\end{equation} \ignorespacesafterend
}
\makeatletter
\newenvironment{itemizedEquation*}{%
	\defineWithinItemizeSpacing%
	\begin{equation*}
}{
	\end{equation*} \ignorespacesafterend
}
\newenvironment{indented}[1][3]{%
	\hfill \begin{minipage}{\dimexpr\textwidth-#1ex} 
	}{
	\end{minipage}
}
\newenvironment{wideitemize}{%
  \begin{itemize}
  \addtolength\itemsep{.4\baselineskip}
}{
  \end{itemize}
}
\newenvironment{narrowItemize}[1][]{%
  \vspace{-.3\baselineskip}%
  \begin{itemize}[#1]
  \addtolength\itemsep{-.1\baselineskip}
}{
  \end{itemize}
}
\newenvironment{narrowEnumerate}[1][1.]{%
  \vspace{-.3\baselineskip}%
  \begin{enumerate}[#1]
  \addtolength\itemsep{-.1\baselineskip}
}{
  \end{enumerate}
}

%% Color definitions
\definecolor{turquoise}{rgb}{0.19, 0.84, 0.78}
\definecolor{mediumturquoise}{rgb}{0.28, 0.82, 0.8}
\definecolor{lava}{rgb}{0.81, 0.06, 0.13}
\colorlet{linkColor}{mediumturquoise}
\colorlet{highlightedTextColor}{lava}

%% Set the color theme of the presentation
\definecolor{jhuBlue}{RGB}{0, 45, 114} % "Heritage" blue; https://brand.jhu.edu/color/
\definecolor{jhuSpiritBlue}{RGB}{114, 172, 229} % lighter blue
\colorlet{themecolor}{jhuBlue}
\colorlet{bgcolor}{themecolor}
\colorlet{textcolor}{white}
\usecolortheme[named=themecolor]{structure}
%\setbeamercolor{titlelike}{fg=black}
\setbeamercolor{frametitle}{fg=themecolor, bg=white}
\setbeamercolor{section in head/foot}{fg=textcolor}
\setbeamercolor{background canvas}{bg=white}
%\colorlet{toccolor}{black}
%\setbeamercolor{section in toc}{fg=toccolor}
\setbeamercolor{block title}{bg=themecolor!75!gray!25}
\setbeamercolor{block body}{bg=themecolor!75!gray!5}


% Macros

%% Utility macros
\newcommand{\noteBullet}{\hspace*{.75em}\textcolor{themecolor}{$\blacktriangleright$}\ }

%% General math macros
\newcommand{\given}{\thinnerspace | \thinnerspace}
\newcommand{\divby}{\thinnerspace /}
\newcommand{\defeq}{\vcentcolon =} % `\coloneqq` and `\vcentcolon` requires the `mathtools` package: https://tex.stackexchange.com/questions/194344/symbol-for-definition
\newcommand{\diff}{\operatorname{\mathrm{d}}\!{}}
\newcommand{\sign}{\mathrm{sign}}
\newcommand{\diag}{\mathrm{diag}}
\newcommand{\medcap}{\mathbin{\mathsmaller{\bigcap}}}
\DeclareMathOperator*{\argmin}{argmin}
\DeclareMathOperator*{\argmax}{argmax}
\newcommand{\elemwiseProd}{\odot}
\renewcommand{\complement}{{\raisebox{1.5pt}{\scriptsize $\mathsf{c}$}}}
\newcommand{\transpose}{\text{\raisebox{.5ex}{$\intercal$}}}
\newcommand{\lowerscript}[1]{\raisebox{-2pt}{\scriptsize $#1$}}
\newcommand{\yesnumber}{\addtocounter{equation}{1}\tag{\theequation}}
\newcommand{\thinnerspace}{\mskip.5\thinmuskip}
\newcommand{\thinnestspace}{\mskip.25\thinmuskip}
\newcommand{\spaceBeforePartial}{\mskip\thinmuskip}
\newcommand{\scriptsumi}{{\scriptsize \sum_i}} % Suscript needs to be within bracket, hence the hard coding

%%% Define a larger "hat" of fixed width: https://tex.stackexchange.com/questions/20473/how-can-i-manually-choose-the-size-of-a-wide-accent-math-mode/20477#20477
\DeclareMathSymbol{\widehatsym}{\mathord}{largesymbols}{"62}
\newcommand{\lowerwidehatsym}{\text{\smash{\raisebox{-1.3ex}{$\widehatsym$}}}}
\newcommand\fixedwidehat[1]{%
  \mathchoice
    {\accentset{\displaystyle\lowerwidehatsym}{#1}}
    {\accentset{\textstyle\lowerwidehatsym}{#1}}
    {\accentset{\scriptstyle\lowerwidehatsym}{#1}}
    {\accentset{\scriptscriptstyle\lowerwidehatsym}{#1}}
}

%% Probability / statistics macros
\newcommand{\probability}{\mathbb{P}}
\newcommand{\expectation}{\mathbb{E}}
\newcommand{\variance}{\mathrm{Var}}
\DeclareMathOperator*{\covariance}{Cov}
\newcommand{\indicator}{\mathds{1}}
\newcommand{\eqDistribution}{\mathrel{\raisebox{-.2ex}{$\overset{\scalebox{.6}{$\, d$}}{=}$}}}
\newcommand{\iidSim}{\mathrel{\raisebox{-.3ex}{$\overset{\text{i.i.d.}}{\sim}$}}}
\newcommand{\unifDist}{\mathrm{Unif}}
\newcommand{\normalDist}{\mathcal{N}}
\newcommand{\betaDist}{\mathrm{Beta}}
\newcommand{\gammaDist}{\mathrm{Gamma}}
\newcommand{\mle}[1]{\widehat{#1}_{\textrm{mle}}}
\newcommand{\map}[1]{\widehat{#1}_{\textrm{map}}}
\newcommand{\empiricalVar}{\widehat{\variance}}
\newcommand{\kldivergence}{D_{\mathrm{KL}}}
\newcommand{\truthSub}{\mathrm{tru}}

%% Variable macros / Aliases
\newcommand{\nPred}{p}
\newcommand{\density}{\pi}
\newcommand{\likelihood}{L}
\newcommand{\infoMat}{\mathcal{I}}
\newcommand{\by}{\bm{y}}
\newcommand{\bx}{\bm{x}}
\newcommand{\bX}{\bm{X}}
\newcommand{\bmu}{\bm{\mu}}
\newcommand{\bbeta}{\bm{\beta}}
\newcommand{\btheta}{\bm{\theta}}
\newcommand{\bPhi}{\bm{\Phi}}
\newcommand{\bSigma}{\bm{\Sigma}}
\newcommand{\minus}{\scalebox{1.0}[1.0]{{\scriptsize -}}}
\newcommand{\Id}{\bm{I}}


\title{\centerline{Bayesian analysis under parametric models}}
\author{%
	Aki Nishimura\\
	Department of Biostatistics%
}
%\institute[]{}
\date{}

\titlegraphic{%
	\begin{picture}(0,0)
		% Try coordinate (155, -150) for a one-line title, (165, -160) if two-lines, and (155, -172) if three-lines
		\put(165, -160){\makebox(0,0)[lb]{\includegraphics[width=.45\textwidth]{jhsph_biostat_logo}}}
	\end{picture}%
}

\begin{document}

\maketitle

%\section{Bayesian estimation of incidence proportion}

\begin{frame}<handout:0>
\frametitle{Example: estimating incidence proportion}
\textbf{Question:} 
\only<+>{``Does the new vaccine cause \ldots?''}%
\only<+->{``Does the new vaccine cause Guillain-Barré syndrome?''}
% Note: Guillain-Barré --- rare but potentially a life-threatening disorder with body's immune system attacking nerves and causing muscle weakness

\begin{minipage}[t][2\baselineskip]{\linewidth}
\only<+>{%
(Unfortunately, this turned out to be the case for the 1976 swine flu vaccine, albeit the syndrome remaining extremely rare.)
}%
\only<+->{%
To answer, we want to estimate the incidence proportion within, say, 4 weeks of vaccination and compare it to the background rate.
}%
\end{minipage}

\pause[\thebeamerpauses]
\textbf{Model:} We assume a binomial model
\begin{tightEquation*}
y_1, \ldots, y_n \given p 
	\iidSim \mathrm{Bernoulli}(p)
\end{tightEquation*}
with the likelihood function
\begin{tightEquation*}
\likelihood(\by \given p) 
	= \textstyle \prod_i p^{y_i} (1 - p)^{1 - y_i}
	= p^{{\scriptscriptstyle \sum_i} y_i} (1 - p)^{n - {\scriptscriptstyle \sum_i} y_i}.
\end{tightEquation*}

\pause
\textbf{Inference:}
Maximizing the likelihood gives us $\mle{p} = {\scriptstyle \sum_i} y_i \divby n$.

\pause 
Problem? 
There's a good chance $\mle{p} = 0$; the null incidence of Guillain-Barré is $\leq 1.5$ in million during 4-week period.

\pause
Potential fix is continuity correction, which yields $\widehat{p} = 0.5 \divby (n + 1)$.
\end{frame}


\begin{frame}<beamer:0>
\frametitle{Example: estimating incidence proportion}
\textbf{Question:} 
\only<+->{``Does the new vaccine cause Guillain-Barré syndrome?''}

\begin{minipage}[t][2\baselineskip]{\linewidth}
\only<+->{%
To answer, we want to estimate the incidence proportion within, say, 4 weeks of vaccination and compare it to the background rate.
}%
\end{minipage}

\pause[\thebeamerpauses]
\textbf{Model:} We assume a binomial model
\begin{tightEquation*}
y_1, \ldots, y_n \given p 
	\iidSim \mathrm{Bernoulli}(p)
\end{tightEquation*}
with the likelihood function
\begin{tightEquation*}
\likelihood(\by \given p) 
	= \textstyle \prod_i p^{y_i} (1 - p)^{1 - y_i}
	= p^{{\scriptscriptstyle \sum_i} y_i} (1 - p)^{n - {\scriptscriptstyle \sum_i} y_i}.
\end{tightEquation*}

\pause
\textbf{Inference:}
Maximizing the likelihood gives us $\mle{p} = {\scriptstyle \sum_i} y_i \divby n$.

\pause 
Problem? 
There's a good chance $\mle{p} = 0$; the null incidence of Guillain-Barré is $\leq 1.5$ in million during 4-week period.

\pause
Potential fix is continuity correction, which yields $\widehat{p} = 0.5 \divby (n + 1)$.
\end{frame}


\begin{frame}
\frametitle{Example: estimating incidence proportion}

Let's try Bayesian analysis of the binomial model with likelihood
\begin{tightEquation*}
\likelihood(\by \given p) 
	= p^{{\scriptscriptstyle \sum_i} y_i} (1 - p)^{n - {\scriptscriptstyle \sum_i} y_i}.
\end{tightEquation*}

\pause
\textbf{Prior:} 
Convenient, reasonably flexible prior here is the beta dist:
\begin{tightEquation*}
\density(p) \propto p^a (1 - p)^b \sim \mathrm{Beta}(a, b)
\end{tightEquation*}
with $\expectation[p] = a \divby (a + b)$ and $\variance[p] = \expectation[p] (1 - \expectation[p]) \divby (a + b + 1)$.

\pause
\vspace*{.25\baselineskip}
\textbf{Posterior:} 
Bayes' rule tells us that
\begin{tightEquation*}
\begin{aligned}
\density(p \given \by)
%	\propto \likelihood(\by \given p) \density(p)
%	&\propto p^{a + \scriptscriptsum_i y_i} (1 - p)^{b + n - \scriptscriptsum_i y_i} \\
	&\sim \mathrm{Beta}(a + {\scriptstyle \sum_i} y_i, b + n - {\scriptstyle \sum_i} y_i).
\end{aligned}
\end{tightEquation*}

\pause
In particular, the posterior mean of $p$ is
\begin{tightEquation*}
\widehat{p}_{\textrm{bayes}} 
	= \expectation[p \given y] 
	= \frac{a + {\scriptstyle \sum_i} y_i}{a + b + n}
	= \frac{(a + b) \thinnerspace \expectation[p] + n \thinnerspace \mle{p}}{a + b + n},
\end{tightEquation*}
\pause
where $a + b$ can be interpreted as \textit{prior sample size}.
% Note: the similarity to MLE. In particular, the posterior mean coincides with MLE in the limit $a, b \to 0$. 
\end{frame}

\begin{frame}
\frametitle{Example: estimating incidence proportion}

Posterior provides much more information than a point estimate. 

\pause
For example, we can calculate a posterior 95\% credible interval by solving for $l$ and $u$ such that
\begin{tightEquation*}
\int_l^u \density(p \given \by) \diff p = 0.95,
\end{tightEquation*}
\pause%
as well as a post predictive dist on future data $\by_{\textrm{new}}$ via relation
\begin{tightEquation*}
\density(\by_{\textrm{new}} \given \by)
	= \int \likelihood(\by_{\textrm{new}} \given p) \thinnerspace \density(p \given \by) \diff p.
\end{tightEquation*}

\pause
\textbf{Note:} Bayesian credible intervals, philosophically speaking, have completely different interpretations from frequentist conf intervals.
\end{frame}

\begin{frame}
\frametitle{Example: estimating incidence proportion}

The remaining question is: how do we choose the prior?

\pause
Bayesians can, for better or worse, incorporate existing scientific knowledge into a prior.

\pause
To properly construct such an informed prior, we have to:
\pause%
\begin{narrowItemize}[<+->]
\item Find every post-vaccine incidence proportion of Guillain-Barré ever reported (and unreported?) in the literature, 
\item Conduct some sort of meta-analysis to obtain a distribution $\density(p)$ for a ``generic'' vaccine, 
\item Somehow account for the difference between mRNA and conventional vaccines, 
%\item Wonder how to incorporate other (e.g.\ pharmacologic) evidence, 
\item Wonder how many other bloody things you have to consider,
\item And maybe start to despair at some point?
\end{narrowItemize}
\end{frame}

\begin{frame}
\frametitle{Example: estimating incidence proportion}

\vspace*{-.3\baselineskip}
As a more pragmatic alternative to constructing a scientifically informed prior, we can consider specifying it in terms of: 
\vspace*{-.1\baselineskip}
\begin{narrowItemize}
\item Prior mean $a \divby (a + b)$,
\item Prior sample size $a + b$, and/or
\item Prior credible intervals.
\end{narrowItemize}
\vspace*{-.1\baselineskip}

\pause
For example, we can 
\vspace*{-.1\baselineskip}
\begin{narrowItemize}
\item Fix the prior sample size to be $1$, and
\item Use a reported background rate of Guillain-Barré ($\approx 1$ per million over 4-weeks) as the prior mean.
\end{narrowItemize}
\vspace*{-.4\baselineskip}
\pause
In other words, solve $(a + b) = 1$ and $a \divby (a + b) = 10^{-6}$ for $a$ and $b$. % which gives $a = 10^{-6}$ and $b = 1 - a$.

\pause
\smallskip
Or, we can choose $a$ and $b$ so that the prior 95\% CI spans the background rate and highest vaccine-induced rate ever reported.%

\end{frame}


\begin{frame}
\frametitle{Example: estimating incidence proportion}

It's also possible to avoid relying on existing evidence by using an \textit{objective prior};
e.g.\ $\betaDist(1/2, 1/2)$ for the binomial model:%

\pause%
\begin{narrowItemize}
\item Removes subjective decisions in Bayesian inference; 
\pause
\item Justified in particular by its frequentist property: Bayesian CI $[q_{\alpha}, 1]$ and $[0, q_{1 - \alpha}]$ achieve coverage at level $1 - \alpha + O(1/n)$;
\pause
\item Excellent frequentist property in non-asymptotic regimes.
\end{narrowItemize}

\pause
The corresponding post mean coincides with the freq estimator based on continuity correction:
\begin{tightEquation*}
\widehat{p}_{\textrm{bayes}} 
	= \expectation[p \given y] 
	= \frac{{\scriptstyle \sum_i} y_i + 1/2}{n + 1}.
\end{tightEquation*}
\end{frame}

\begin{frame}
\frametitle{Example: estimating incidence proportion}

The prior $\propto p^{-1/2} (1 - p)^{-1/2}$ may not be the most obvious choice.
\onslide<2->{But the more obvious ``flat'' prior has its own problems.}

\vspace{.15\baselineskip}
\begin{minipage}{.48\linewidth}
	\begin{figure}
		\begin{tikzpicture}
		\begin{axis}[
			width=0.575\paperheight,
			height=0.575\paperheight,
		    axis lines = left,
		    xmin=0,
		    xmax=1,
		    ymin=0,
		    ymax=3.5,
		    xlabel = \(p\),
		    ylabel = {\(\)},
		    legend style={
		    	font=\small, 
		    	draw=none, % Remove legend box
		    	empty legend % Remove line marker
		    }
		]
		%Below the red parabola is defined
		\addplot [
			domain=0.01:0.99,
		    samples=100, 
		    color=jhuBlue,
		    thin
		]
		{x^(-1/2) * (1 - x)^(-1/2) / pi};
		\addlegendentry{\({\tiny \betaDist\!\left(\thinnerspace p \thinnerspace \middle| \thinnerspace \frac12, \frac12\right)}\)}
		\end{axis}
		\end{tikzpicture}
	\end{figure}
\end{minipage}
~
\onslide<2->{%
\begin{minipage}{.48\linewidth}
	\begin{figure}
		\begin{tikzpicture}
		\begin{axis}[
			width=0.575\paperheight,
			height=0.575\paperheight,
		    axis lines = left,
		    xmin=0,
		    xmax=1,
		    ymin=0,
		    ymax=3.5,
		    xlabel = \(p\),
		    ylabel = {\(\)},
		    legend style={
		    	font=\small, 
		    	draw=none, % Remove legend box
		    	empty legend % Remove line marker
		    }
		]
		%Below the red parabola is defined
		\addplot [
			domain=0.01:0.99,
		    samples=100, 
		    color=jhuSpiritBlue,
		    thin
		]{1};
		\addlegendentry{\({\tiny \betaDist\!\left(\thinnerspace p \thinnerspace \middle| \thinnerspace 1, 1\right)}\)}
		\end{axis}
		\end{tikzpicture}
	\end{figure}
\end{minipage}
}%
\vspace{.15\baselineskip}
\onslide<3>{%
At the least, Bayesian inference has an explicit, clear interpretation in terms of the prior.
}
\end{frame}


\begin{frame}
\frametitle{Homework}
\textbf{Problem 1:} 
The Jeffreys prior is defined as $\density(\btheta) = \det\!\left(\infoMat(\btheta) \right)^{1/2}$, where $\infoMat(\btheta)$ is the Fisher information, and is arguably the prior of choice for objective Bayesian analysis when $\btheta$ is univariate.
Show that $p \sim \betaDist(\frac12, \frac12)$ is the Jeffreys prior for the bernoulli model.
What would be the Jeffreys prior if we use an alternate parametrization $\log \frac{p}{1 - p}$? 
How does that affect posterior inference?
\end{frame}


\begin{frame}
\frametitle{Homework}
\textbf{Problem 2:} 
Consider posterior inference for binomial data $y_1, \ldots, y_{25} \given p \iidSim \mathrm{Bernoulli}(p)$ under the prior $p \sim \betaDist(\frac12, \frac12)$.
Numerically study the frequentist coverage of Bayesian 95\% credible intervals under $p = \frac{k}{100}$ for $k = 1, \ldots, 99$.
You may analytically/numerically calculate any necessary quantities or resort to Monte Carlo simulation.
Compare the credible intervals' performance to the coverage achieved by a frequentist procedure, such as the one based on the asymptotic normality of {\small MLE} or any other techniques that may be better suited to this specific setting.
\end{frame}


\section{\centerline{Gaussian likelihood models}}


\begin{frame}
\frametitle{Univariate case}

Consider the problem of estimating a population mean $\mu$ under
\begin{tightEquation*}
y_1, \ldots, y_n \given \mu, \sigma \iidSim \normalDist(\mu, \sigma^2).
\end{tightEquation*}

\pause
For simplicity, let's first consider $\sigma$ to be fixed.
Also, for convenience, we parametrize in terms of \textit{precision} $\phi = \sigma^{-2}$.
\end{frame}

\newcommand{\priorPrec}{\phi_0}
\newcommand{\postPrec}{\phi_{\textrm{post}}}
\begin{frame}
\frametitle{Univariate case}
Convenient prior on $\mu$ is Gaussian:\ under the prior $\mu \sim \normalDist(\mu_0, \priorPrec^{-1})$,%
\begin{tightEquation*}
\begin{aligned}
\density(\mu \given \bm{y}, \phi)
	&\propto \likelihood(\bm{y} \given \mu, \phi ) \density(\mu) \\
	\pause
%	&\propto \exp\!\left( - \frac{\phi}{2} \sum_i (y_i - \mu)^2\right)
%		\exp\!\left( - \frac{\priorPrec}{2} (\mu - \mu_0)^2 \right) \\
	&\propto \phi^{n / 2} \exp\!\left( - \frac12 \left[
			\phi \sum_i (y_i - \mu)^2
			+ \priorPrec (\mu - \mu_0)^2
		\right] \right) \\
	\pause
	&\propto \exp\!\left( - \frac12 \left[
			\phi \left( n \mu^2 - 2 \mu {\textstyle \sum_i} y_i \right)
			+ \priorPrec (\mu^2 - 2 \mu \mu_0)
		\right] \right) \\
	\pause
	&= \exp\!\left( - \frac12 \left[
			(n \phi + \priorPrec) \mu^2 
			- 2 \mu \big( \phi {\textstyle \sum_i} y_i + \priorPrec \mu_0 \big)
		\right] \right) \\
	\pause
	&= \exp\!\left( - \frac12 \left[
			\postPrec \left(
				\mu^2 - 2 \mu \mu_{\textrm{post}} 
			\right) 
		\right] \right) 
\end{aligned}
\end{tightEquation*}
for $\postPrec = n \phi + \priorPrec$ and $\mu_{\textrm{post}} = \postPrec^{-1} \left( \phi {\textstyle \sum_i} y_i + \priorPrec \mu_0 \right)$.

\pause
\vspace*{.3\baselineskip}
In other words, the posterior is:
$
\mu \given \bm{y}, \phi
	\sim \normalDist\!\left( \mu_{\textrm{post}}, \postPrec \right).
$
\end{frame}

\begin{frame}
\frametitle{Univariate case}
Note that, as in the beta-binomial case, the posterior mean is a weighted average of prior mean and {\small MLE}:
\begin{tightEquation*}
\mu_{\textrm{post}} = \frac{1}{n \phi + \priorPrec} \left( n \phi \mle{\mu} + \priorPrec \mu_0 \right)
\ \text{ where } \, \mle{\mu} = {\textstyle \sum_i} y_i / n.
\end{tightEquation*}

\pause
Objective ``prior'' for this model is a flat dist $\density(\mu) \propto 1$, which can be though of as the limit $\priorPrec \to 0$ of diminishing prior info. 
%\hspace*{.5em}\textcolor{themecolor}{$\blacktriangleright$} 

\pause
The flat prior is an example of \textit{improper priors}, whose use is justified by the posterior being a proper distribution.

\pause
Under the flat prior, the posterior becomes 
$\mu \given \bm{y}, \phi \sim \normalDist\!\left( n \phi, \mle{\mu} \right)$,
inferences under which coincides with frequentist ones.

\pause
\textbf{Note:} The flat prior (and $\betaDist(1/2, 1/2)$) also has a formal justification as being a Jeffreys/reference prior.
\end{frame}

\begin{frame}
\frametitle{Univariate case: side note on conjugate priors}
We've seen that a Gaussian (beta) prior leads to a Gaussian (beta) posterior under the Gaussian (binomial) likelihood model.

\pause
\smallskip
More generally, a prior is said to be \textit{conjugate} for the likelihood if the corresponding posterior belongs to the same distribution family.

\pause
\smallskip
Conjugate priors lead to (semi) analytically tractable posteriors and%
\begin{narrowItemize}[<+->]
\item play an important, if not as critical as in the pre-Stan era, role in computationally tractable Bayesian modeling;
\item provide insights into the mechanics of Bayesian inference.
\end{narrowItemize}

\end{frame}

\begin{frame}
\frametitle{Univariate case with unknown variance}
\pause
Easy enough to infer by examining the likehood
\begin{tightEquation*}
\likelihood(\bm{y} \given \mu, \phi )
	\propto \phi^{n/2} \exp\!\left( - \frac12 
			\phi \sum_i (y_i - \mu)^2
		\right)
\end{tightEquation*}
\pause%
that a (conditionally) conjugate prior for $\phi$ is a gamma distribution:
\begin{tightEquation*}
\density(\phi) 
	\propto \phi^{a - 1} \exp(-b \phi)
	\sim \gammaDist(\textrm{shape} = a, \thinnerspace \textrm{rate} = b)
\end{tightEquation*}
with $\expectation[\phi] = a / b$ and $\variance[\phi] = a / b^2$.

\pause 
\smallskip
The corresponding posterior is
\begin{tightEquation*}
\phi \given \bm{y}, \mu 
	\sim \gammaDist\!\left(a + \frac{n}{2}, \, b + \frac{1}{2} \textstyle \sum_i (y_i - \mu)^2 \right).
\end{tightEquation*}
\end{frame}

\begin{frame}
\frametitle{Univariate case with unknown variance}
For interpretability, more helpful to parametrize the prior as 
\begin{tightEquation*}
\phi \sim \gammaDist\!\left( \frac{\nu_0}{2}, \thinnerspace \frac{\nu_0 \varphi_0^{-1}}{2} \right),
%\gammaDist\!\left( \textrm{shape} = \frac{\nu_0 \varphi_0}{2}, \thinnerspace \textrm{rate} = \frac{\nu_0}{2} \right)
\end{tightEquation*} 
under which $\expectation[\phi] = \varphi_0$ and $\variance[\phi] = 2 \varphi_0 / \nu_0$.

\medskip
The posterior now takes the form
\begin{tightEquation*}
\phi \given \bm{y}, \mu 
	\sim \gammaDist\!\left(
		\frac{\nu_0 + n}{2}, \, % \frac{1}{2} (\nu_0 + n), \, 
		\frac{\nu_0 \varphi_0^{-1} + \textstyle \sum_i (y_i - \mu)^2}{2} % \frac{1}{2} \left(\nu_0 \varphi_0^{-1} + \textstyle \sum_i (y_i - \mu)^2 \right) 
	\right)
\end{tightEquation*}
with
\begin{tightEquation*}
\expectation\left[ \phi \given \bm{y}, \mu \right]^{-1}
	= \frac{1}{\nu_0 + n} \left( \nu_0 \varphi_0^{-1} + n \thinnerspace \frac{\textstyle \sum_i (y_i - \mu)^2}{n} \right).
\end{tightEquation*}

\pause
\smallskip
In particular, $\nu_0$ can be interpreted as a prior degrees of freedom.
\end{frame}

\begin{frame}
\frametitle{Univariate case with unknown variance}
Objective prior $\density(\phi) \propto 1 / \phi$ coincides with the limit $\nu_0 \to 0$ and yields the posterior
\begin{tightEquation*}
\phi \given \bm{y}, \mu 
	\sim \gammaDist\!\left(
		\frac{n}{2}, \,
		\frac{\textstyle \sum_i (y_i - \mu)^2}{2} % \frac{1}{2} \left(\nu_0 \varphi_0^{-1} + \textstyle \sum_i (y_i - \mu)^2 \right) 
	\right)
\end{tightEquation*}
with
\begin{tightEquation*}
\expectation\left[ \phi \given \bm{y}, \mu \right]^{-1}
	= \frac{\textstyle \sum_i (y_i - \mu)^2}{n}
	= 1 \divby \mle{\phi}.
\end{tightEquation*}

\pause
Posterior becomes dominated by likelihood as prior becomes more diffuse, so the agreement with {\small MLE} isn't entirely surprising but \ldots

\pause
it's also worth noting the posterior mean differs from the mode
\begin{tightEquation*}
{\textstyle \argmax_{\phi}} \density\!\left( \phi \given \bm{y}, \mu \right)
	= \max\!\left\{
		0, \, \frac{n - 2}{\textstyle \sum_i (y_i - \mu)^2}
	\right\}.
\end{tightEquation*}
\end{frame}


\begin{frame}
\frametitle{Univariate case with unknown variance}
We now turn to the problem of estimating $\mu$ and $\phi$ simultaneously.

\pause
Most obvious choice is to assume $\mu$ and $\phi$ are a priori independent:
\begin{tightEquation*}
\density(\mu, \phi) = 
	\normalDist\!\left( \mu \given \mu_0, \priorPrec \right)
	\times
	\gammaDist\!\left( \phi \given \nu_0 / 2, \thinnerspace \nu_0 \varphi_0^{-1} / 2 \right).
\end{tightEquation*}

\pause
This prior, however, is only conditionally conjugate; i.e.\ 
while $\mu \given \bm{y}, \phi$ and $\phi \given \bm{y}, \mu$ belong to the original families, $\mu, \phi \given \bm{y}$ does not.

\pause
\textbf{Note:} The independent prior is otherwise a perfectly good choice!

Fully conjugate priors are worthy additions to your toolbox when they exist, but there is no reason to be religious about them.
\end{frame}


\begin{frame}
\frametitle{Univariate case with unknown variance}
To identify a conjugate prior, we let $\bar{y} = n^{-1} \sum_i y_i$ and observe 
\begin{tightEquation*}
\begin{aligned}
\likelihood(\bm{y} \given \mu, \phi )
	&\propto \phi^{n/2} \exp\!\left( - \frac12 
			\phi \sum_i (y_i - \mu)^2
		\right) \\
	&\propto \phi^{n/2} \exp\!\left( - \frac12 
			\phi \sum_i (y_i - \bar{y})^2
		\right)
		\exp\!\left( - \frac12 
			n \phi (\mu - \bar{y})^2
		\right)
\end{aligned}
\end{tightEquation*}
b/c $\sum_i (y_i - \mu)^2 = \sum_i (y_i - \bar{y})^2 + \sum_i (\bar{y} - \mu)^2$.

\pause
\medskip
After staring at the last expression long enough, you'd realize that
\begin{tightEquation*}
\likelihood(\bm{y} \given \mu, \phi )
	\propto 
		\gammaDist\!\left( \phi \, \middle| \,  \frac{n + 1}{2}, \, \frac{\sum_i (y_i - \bar{y})^2}{2} \right)  
		\normalDist\!\left( \mu \, \middle| \, \bar{y}, \, \frac{1}{n \phi} \right).
\end{tightEquation*}
% Note: the "-1" in $n - 1$ comes from the need for $\phi^{n/2}$ for the Gaussian density.
\end{frame}


\begin{frame}
\frametitle{Univariate case with unknown variance}
It's now easy enough to guess the following prior as conjugate:
\begin{tightEquation*}
\phi \sim \gammaDist\!\left( \frac{\nu_0}{2}, \thinnerspace \frac{\nu_0 \varphi_0^{-1}}{2} \right), \
\mu \given \phi \sim \normalDist\!\left( \mu_0, \thinnerspace \frac{1}{\kappa_0 \phi} \right)
\vspace*{-.15\baselineskip}% 
\end{tightEquation*}
with density 
\vspace*{-.15\baselineskip}
\begin{tightEquation*}
\density(\mu, \phi)
	\propto \phi^{\frac{\nu_0}{2} - 1} \exp\!\left( - \frac{\nu_0 \varphi_0^{-1}}{2} \phi \right)
		\exp\!\left( - \frac{1}{2} \kappa_0 \phi (\mu - \mu_0)^2 \right).
\end{tightEquation*}

\pause
\medskip
The corresponding posterior is given by (homework)
\begin{tightEquation*}
\begin{aligned}
&\phi \given \bm{y} 
	\sim \gammaDist\!\left(
		\frac{\nu_0 + n}{2}, \thinnerspace b_{\textrm{post}}%
	\right), \\
&\mu \given \bm{y}, \phi 
	\sim \normalDist\!\left( \frac{\kappa_0 \mu_0 + n \bar{y}}{\kappa_0 + n}, \thinnerspace \frac{1}{(\kappa_0 + n) \phi} \right),
\end{aligned}
\end{tightEquation*}
where 
$\, b_{\textrm{post}} = {\displaystyle \frac12} \left(
	\nu_0 \varphi_0^{-1} 
	+ \textstyle \sum_i (y_i - \bar{y})^2
	+ \frac{\kappa_0 n}{\kappa_0 + n} (\mu_0 - \bar{y} )^2
\right)$.
\end{frame}


\begin{frame}
\frametitle{Univariate case with unknown variance}
The posterior estimates mostly have the same structure as before, but let's take a look at
\begin{equation*}
\expectation[ \phi \given \bm{y} ]
	= \frac{1}{\nu_0 + n} \left(
		\nu_0 \varphi_0^{-1} 
		+ \textstyle \sum_i (y_i - \bar{y})^2
		+ \frac{\kappa_0 n}{\kappa_0 + n} (\mu_0 - \bar{y} )^2
	\right).
\end{equation*}

\pause
The $\frac{\kappa_0 n}{\kappa_0 + n} (\mu_0 - \bar{y} )^2$ part shows one way in which $\mu$ and $\phi$ interact:
\begin{tightEquation*}
\text{%
	discrepancy between $\mu_0$ and $\bar{y}$ % Note: $\mu_0$ = prior guess and $\bar{y}$ = empirical mean 
	$\, \Rightarrow \,$ inflated estimate of $\phi$.
}%
\end{tightEquation*}

\pause
\smallskip
But note also this ``bias'' diminishes as $n \to \infty$ since $\frac{\kappa_0 n}{\kappa_0 + n} \leq \kappa_0$.
\end{frame}


\begin{frame}
\frametitle{Univariate case with unknown variance}
The prior has a feature that the apriori (conditional) variance $\variance(\mu \given \phi) = 1 / \kappa_0 \phi$ depends on the noise variance $1 / \phi$.

\pause
\smallskip
The prior thus has the behavior ``smaller the $\phi$ and noisier the $y_i$'s, less informative the prior on $\mu \given \phi$.''

\pause
\smallskip
This may not be an accurate reflection of the prior knowledge, but is a reasonable data-adaptive behavior.
\end{frame}


\begin{frame}
\frametitle{Univariate case with unknown variance}
Identifying an appropriate objective prior isn't always straight- forward in multi-parameter settings, but the established one here is:
\begin{tightEquation*}
\density(\mu, \phi) \propto 1 / \phi.
\end{tightEquation*}

\pause
The objective prior yields a posterior
\begin{tightEquation*}
%\begin{aligned}
\phi \given \bm{y} 
	\sim \gammaDist\!\left(
		\frac{n - 1}{2},
		\frac{\textstyle \sum_i (y_i - \bar{y})^2}{2} 
	\right), \
\mu \given \bm{y}, \phi 
	\sim \normalDist\!\left( \bar{y}, \thinnerspace \frac{1}{n \phi} \right)
%\end{aligned}
\end{tightEquation*}
with
\begin{tightEquation*}
\expectation\!\left[ 
	(\mu, \phi) \given \bm{y}
\right]
	= \left( 
		\bar{y}, \, \frac{n - 1}{\textstyle \sum_i (y_i - \bar{y})^2}
	\right).
\end{tightEquation*}
\end{frame}


\section{\centerline{Linear regression}}


\begin{frame}
\frametitle{Linear model with Gaussian errors}
Consider a regression model with design matrix $\bX \in \mathbb{R}^{n \times \nPred}$
\begin{tightEquation*}
\bm{y} = \bX \bbeta + \bm{\epsilon} 
	\ \text{ for } \,
	\bm{\epsilon} \sim \normalDist\!\left(\bm{0}, \phi^{-1} \bm{I} \right);
\end{tightEquation*}
or, to be more precise, 
\begin{tightEquation*}
\bm{y} \given \bX, \bbeta, \phi
	\sim \normalDist\!\left( \bX \bbeta, \phi^{- 1} \bm{I} \right).
\end{tightEquation*}

\pause
As we will see, posterior inference under this linear model is structurally very similar to the univariate Gaussian model.
\end{frame}

\begin{frame}
\frametitle{Prior choice}
The linear model has the likelihood
\begin{tightEquation*}
\begin{aligned}
\likelihood(\bm{y} \given \bX, \bbeta, \phi)
	&\propto \phi^{n / 2} \exp\!\left(
		- \frac{\phi}{2} \left\| \bm{y} - \bX \bbeta \right\|^2 
	\right) \\
	&\propto \phi^{n / 2} \exp\!\left(
		- \frac{1}{2} \left( 
			\phi \thinnerspace \bbeta^\transpose \bX^\transpose \bX \bbeta
			- 2 \thinnerspace \phi \thinnerspace \bbeta^\transpose \bX^\transpose \bm{y}
		\right) 
	\right).
\end{aligned}
\end{tightEquation*}

\pause
As a function of $\bbeta$, the likelihood looks like a Gaussian density
\begin{tightEquation*}
\normalDist\!\left( 
	\thinnerspace \bbeta \, \thinnerspace \middle| \thinnerspace
	\left( \bX^\transpose \bX \right)^{-1} \bX^\transpose \bm{y},
	\left( \phi \thinnerspace \bX^\transpose \bX \right)^{-1} 
\right).
\end{tightEquation*}

\pause
As a function of $\phi$,  the likelihood looks like a Gamma density.
%\begin{tightEquation*}
%\gammaDist\!\left(
%	\thinnerspace \phi \, \thinnerspace \middle| \thinnerspace
%	\frac{n}{2} + 1, 
%	\frac{\left\| \bm{y} - \bX \bbeta \right\|^2}{2}
%\right).
%\end{tightEquation*}

\pause
\smallskip
So you can guess the following prior to be conditionally conjugate:
\begin{tightEquation*}
\bbeta \sim \normalDist(\bmu_0, \bPhi_0^{-1})
\ \text{ and } \
\phi \sim \gammaDist\!\left( \frac{\nu_0}{2}, \thinnerspace \frac{\nu_0 \varphi_0^{-1}}{2} \right).
\end{tightEquation*}
\end{frame}


\begin{frame}
\frametitle{Posterior under (conditionally) conjuage priors}
In particular, note that $\bbeta \sim \normalDist(\bmu_0, \bPhi_0)$ means
\begin{tightEquation*}
\begin{aligned}
\density(\bbeta)
%	&\propto \exp\!\left(
%		- \frac{1}{2} (\bbeta - \bmu_0)^\transpose \bPhi_0 (\bbeta - \bmu_0)
%	\right) \\
	&\propto \exp\!\left(
		- \frac{1}{2} \left( 
			\bbeta^\transpose \bPhi_0 \bbeta 
			- 2 \bbeta^\transpose \bPhi_0 \bmu_0
		\right) 
	\right).
\end{aligned}
\end{tightEquation*}
\pause
So we have
\begin{tightEquation*}
\begin{aligned}
\density(\bbeta \given \bm{y}, \bX, \phi)
	&\propto \likelihood(\bm{y} \given \bbeta, \bX, \phi) \thinnerspace \density(\bbeta) \\
	&\propto \exp\!\left(
		- \frac{1}{2} \left( 
			\phi \thinnerspace \bbeta^\transpose \bX^\transpose \bX \bbeta
			- 2 \thinnerspace \phi \thinnerspace \bbeta^\transpose \bX^\transpose \bm{y}
		\right)
	\right) \\
	&\hphantom{\propto}\ \exp\!\left(
		- \frac{1}{2} \left( 
			\bbeta^\transpose \bPhi_0 \bbeta 
			- 2 \bbeta^\transpose \bPhi_0 \bmu_0
		\right) 
	\right) \\
	&\propto \exp\!\left(
		- \frac{\phi}{2} \left( 
			\bbeta^\transpose \bPhi_{\textrm{post}} \bbeta
			- 2 \bbeta^\transpose \bPhi_{\textrm{post}} \bmu_{\textrm{post}}
		\right) 
	\right),
\end{aligned}
\end{tightEquation*}
where $\bPhi_{\textrm{post}} = \phi \thinnerspace \bX^\transpose \bX + \bPhi_0$ and $\bmu_{\textrm{post}} = \bPhi_{\textrm{post}}^{-1} \left( \phi \thinnerspace \bX^\transpose \bm{y} + \bPhi_0 \bmu_0 \right)$.
\end{frame}


\begin{frame}
\frametitle{Posterior under (conditionally) conjuage priors}

In other words, the posterior conditional for $\bbeta$ is:
\begin{tightEquation*}
\bbeta \given \bm{y}, \bX, \phi
	\sim \normalDist\!\left(
		 \bmu_{\textrm{post}},
		 \bPhi_{\textrm{post}}
	\right)
\end{tightEquation*}
for $\bmu_{\textrm{post}} = \bPhi_{\textrm{post}}^{-1} \left( \phi \thinnerspace \bX^\transpose \bm{y} + \bPhi_0 \bmu_0 \right)$ and $\bPhi_{\textrm{post}} = \phi \thinnerspace \bX^\transpose \bX + \bPhi_0$.

\pause
\medskip
We similarly obtain the posterior conditional for $\phi$:
\begin{tightEquation*}
\phi \given \bm{y}, \bX, \bbeta 
	\sim \gammaDist\!\left( 
		\frac{\nu_0 + n}{2}, \thinnerspace \frac{\nu_0 \varphi_0^{-1} + \| \bm{y} - \bX \bbeta \|^2}{2} 
	\right).
\end{tightEquation*}

\pause
As in the univariate case, we could also consider a conjugate prior
\begin{tightEquation*}
\bbeta \given \phi \sim \normalDist(\bmu_0, (\kappa_0 \phi)^{-1} \bm{I})
\ \text{ and } \
\phi \sim \gammaDist\!\left( 
		\frac{\nu_0}{2}, \thinnerspace \frac{\nu_0 \varphi_0^{-1}}{2} 
	\right).
\end{tightEquation*}
% Homework idea: derive the corresponding posterior.
% Homework idea: derive the posterior under non iid noise.
\end{frame}


\begin{frame}
\frametitle{Objective Bayesian analysis in linear regression}
Standard objective prior here is, unsurprisingly, $\thinnerspace \density(\bbeta, \phi) \propto 1 / \phi$.

\pause
The corresponding posterior is
\begin{tightEquation*}
\begin{gathered}
\bbeta \given \bm{y}, \bX, \phi
	\sim \normalDist\!\left(
		 \left(  \bX^\transpose \bX \right)^{-1} \! \bX^\transpose \bm{y}, \,
		\left( \phi \bX^\transpose \bX \right)^{-1}
	\right), \\
\phi \given \bm{y}, \bX
	\sim \gammaDist\!\left( 
		\frac{n - \nPred}{2}, \thinnerspace \frac{\| \bm{y} - \bX \mle{\bbeta} \|^2}{2} 
	\right). 
	% Note: $\mle{\bbeta} = \left( \bX^\transpose \bX \right)^{-1} \! \bX^\transpose \bm{y}$
\end{gathered}
\end{tightEquation*}

\pause
\textbf{Note:} The posterior is proper only if $n > \nPred$:
\begin{narrowItemize}[<+->]
\item Flat prior in $\mathbb{R}^\nPred$ is even ``more improper/diffuse'' than in $\mathbb{R}$ because prior uncertainties in the parameters add up.
\item More generally, using improper or proper-but-diffuse priors is a risky business in high-dimensional (i.e.\ large $\nPred$) applications.
\item Requirement $n > \nPred$ is intuitive enough here but it's less obvious in more complex models (where Bayes really shines).
\end{narrowItemize}
\end{frame}


\begin{frame}
\frametitle{Informal objective-ish prior in linear regression}
Let's try constructing a proper yet arguably conservative prior for the linear model.

\pause
(``Conservative'' means biasing, if any, the inference towards the null and encoding less prior information than we actually have.)

\pause
To help us achieve this goal, let's first pre-process the data:
\begin{narrowEnumerate}
\item Center $\bm{y}$ and $\bx_j$'s. % all (both cont.\ and binary) $\bx_j$'s.
%  \item For each binary $\bx_j$, set the more freq category as its baseline; % (i.e. assigning it the value of $0$)
%  this ensures $x_{ij} = 0$ to better represent an average data point.
\item Scale $\bm{y}$ and cont.\ $\bx_j\mkern-1mu$'s so that $\widehat{\variance}(\by) = 1$ and $\widehat{\variance}(\bx_j \mkern-2mu) = 1/4$.%
\end{narrowEnumerate}

\pause
\textbf{Note:} $\widehat{\variance}(\bx) = 1/4$ follows \citet{gelman2008default_prior} and is an attempt to put cont.\ and binary predictors on the same scale. 
\end{frame}


\begin{frame}
\frametitle{Informal objective-ish prior in linear regression}
With the pre-processed data $(\by, \bX)$, we can argue $| \beta_j | = 2$ to represent a reasonable upper bound on the effect sizes because:
\begin{tightEquation*}
| \beta_j | = 2
	\ \Rightarrow \ \widehat{\variance}(\beta_j \bx_j) = \widehat{\variance}(\by).
	% Note: In other words, the variation in the $
\end{tightEquation*}
In other words, the variation in $\bx_j$ explains most of that in $\by$. 

\pause
The prior choice $\beta_j \sim \normalDist(0, 1)$ hence seems reasonable. % with most of its probability in the range $[-2, 2]$.

\pause
\smallskip
As for the intercept, it seems reasonable to set $\beta_0 \sim \normalDist(0, 1)$ so that its prior variance matches that of $\by$.

\end{frame}

\begin{frame}
\frametitle{Informal objective-ish prior in linear regression}
For the residual variance $\phi^{-1}$, we can expect $\phi^{-1} \lesssim 1 = \widehat{\variance}(\by)$.

\pause
If we were to do without conjugacy, we might choose a prior like 
\begin{tightEquation*}
\sigma = \phi^{-1/2} \sim \normalDist^+(0, 1).
	% Note: the prior variance of 1 / 4 would favor \sigma < 1, which isn't ideal since \sigma may well be very close to 1. 
	\vspace*{-.1\baselineskip}%
\end{tightEquation*}
(Also possible to incorporate a soft lower bound, say, $\sigma \geq 0.1$.)

\pause
\smallskip
Gamma prior on $\phi$ is less flexibility but, based the interpretation of $\nu_0$ as the prior degrees of freedom,
	% Note: In particular, the right exponential tail necessarily penalizes larger $\phi$ and the left polynomial tail cannot aggressively penalize small $\phi$.
\begin{tightEquation*}
\phi \sim \gammaDist(\nu_0 / 2, \nu_0 / 2)
	\ \text{ for } \, \nu_0 \in [0.1, \thinnerspace 1]
\end{tightEquation*}
with $\expectation[\phi] = 1 = \widehat{\variance}(\by)$ is defensible.
% Note: BUGS uses $\nu_0 = 0.002$.

\pause
\smallskip
\textbf{Note:} 
Typically, as long as the prior is in a reasonable range, the inference isn't too sensitive.
So don't fret about gory details.
% Note: Plus, 1) you can always do a sensitivity analysis; and 2) if the inference is sensitive, most likely it means you don't have enough data to avoid informative priors.
\end{frame}


\begin{frame}
\frametitle{Informal objective-ish prior in linear regression}
Incidentally, scaling $\by$ and $\bx_j$'s is equivalent to scaling the priors.

\pause
\smallskip
In fact, had we only centered them w/o scaling, the same inference would be obtained under the following prior:
\begin{tightEquation*}
\begin{gathered}
\beta_{\mathrm{cont}, \thinnerspace j} \sim\normalDist\!\left( 0, \, \frac{ \empiricalVar(\by) }{ 4 \empiricalVar(\bx_j) } \right), \ \,
\beta_{\mathrm{bin}, \thinnerspace j} \sim\normalDist\!\left( 0, \, \frac{ \empiricalVar(\by) }{ \empiricalVar(\bx_j) } \right), \\
\beta_0 \sim \normalDist\!\left( 0, \empiricalVar(\by) \right), \ \,
\phi \sim \gammaDist\!\left(
	\frac{\nu_0}{2}, 
	\frac{\nu_0}{2} \thinnerspace \empiricalVar(\by)^{-1}
\right).
%\beta_j \sim\normalDist\!\left( 0, \, \empiricalVar(\by)^{1/2} \big/ \, \empiricalVar(\bx_j)^{1/2} \right) \\
\end{gathered}
\end{tightEquation*}

\pause
\smallskip
Centering variables interacts with the inference as well; e.g.\ had we decided to center $\bx_j$'s but not $\by$, the equivalent inference requires
\begin{tightEquation*}
\beta_0 \sim \normalDist\!\left( \bar{y}, \thinnerspace \empiricalVar(\by) \right).
\end{tightEquation*}
\end{frame}


\begin{frame}
\frametitle{Homework: Gaussian likelihood models}
First problem is optional for those who have taken ``Advanced Methods in Biostatistics'' and the second optional for the rest.

\textbf{Problem 1:} Verify that the posterior under the model $y_i \given \mu, \phi \sim \normalDist(\mu, \phi^{-1})$ and conjugate prior 
\begin{tightEquation*}
\phi \sim \gammaDist\!\left( \frac{\nu_0}{2}, \thinnerspace \frac{\nu_0 \varphi_0^{-1}}{2} \right), \
\mu \given \phi \sim \normalDist\!\left( \mu_0, \thinnerspace \frac{1}{\kappa_0 \phi} \right)
\end{tightEquation*}
is given by 
\begin{tightEquation*}
\begin{aligned}
&\phi \given \bm{y} 
	\sim \gammaDist\!\left(
		\frac{\nu_0 + n}{2}, \thinnerspace b_{\textrm{post}}%
	\right), \\
&\mu \given \bm{y}, \phi 
	\sim \normalDist\!\left( \frac{\kappa_0 \mu_0 + n \bar{y}}{\kappa_0 + n}, \thinnerspace \frac{1}{(\kappa_0 + n) \phi} \right),
\end{aligned}
\end{tightEquation*}
where 
$\, b_{\textrm{post}} = {\displaystyle \frac12} \left(
	\nu_0 \varphi_0^{-1} 
	+ \textstyle \sum_i (y_i - \bar{y})^2
	+ \frac{\kappa_0 n}{\kappa_0 + n} (\mu_0 - \bar{y} )^2
\right)$.
\end{frame}


\begin{frame}
\frametitle{Homework: Gaussian likelihood models}
\textbf{Problem 2:}
%In fitting a regression model to an outcome $\by$ and predictors $\bx_j$, consider scaling 
Consider fitting the linear regression model to a scaled outcome $\by_{\textrm{scl}} = \gamma_Y^{-1} \by$ and scaled predictors $\bx_{\mathrm{scl}, \thinnerspace j} = \gamma_j^{-1} \bx_j$ under a prior $\bbeta \sim \normalDist(\bmu_0, \bSigma_0)$ and $\phi \sim \gammaDist(a, b)$. 
If we were to keep the outcome and predictors in the original scale, how must we modify the prior In order to obtain an equivalent inference?
In what sense(s) the inference ``equivalent''?

\hfill (The problem continues to the next page.)
\end{frame}


\begin{frame}
\frametitle{Homework: Gaussian likelihood models}
\textbf{Problem 2 (cont'd):}
Next, we study the effect of centering variables. 
For this purpose, as you will find out, it is useful to separate out the intercept term from the design matrix and consider the model $\by = \beta_0 \bm{1} + \bX \bbeta + \bm{\epsilon}$ where $\bm{1}$ denotes a vector of $1$'s.
Let $\by_{\textrm{ctr}} = \by - \bar{y} \thinnerspace \bm{1}$ and $\bx_{\mathrm{ctr}, \thinnerspace j} = \bx_j - \bar{x}_j \bm{1}$, where $\bar{y} = n^{-1} \sum_i y_i$ and$\bar{x}_j = n^{-1} \sum_i x_{ij}$.
%In this part of the problem, it is convenient to separate out the intercept $\beta_0$ from the rest of coefficients $\bbeta$.
Consider fitting the regression model with $\by_\mathrm{ctr}$ and $\bX_\mathrm{ctr}$ (i.e.\ $\by_\mathrm{ctr} = \beta_0 \bm{1} + \bX_\mathrm{ctr} \bbeta + \bm{\epsilon}$) under a prior $\beta_0 \sim \normalDist(0, 1)$, $\bbeta \sim \normalDist(\bm{0}, \Id)$, and $\phi \sim \gammaDist(a, b)$.
If we were to fit the regression model to the original variables $\by$ and $\bX$, how must we modify the prior to obtain an equivalent inference?
\end{frame}


\section{\centerline{Asymptotic properties of Bayesian inference}}

\begin{frame}
\frametitle{Posteriors' asymptotic behavior}
We've seen examples in which posteriors yield comparable/ identical inferences as established frequentist procedures.

\pause
In fact, under regular parametric models, Bayesian inference is \textit{guaranteed} to coincide with a frequentist one asymptotically.

\pause
In a sense, this is to be expected from the fact 
\begin{tightEquation*}
\log \density(\btheta \given \by^n)
	= \mathrm{const} + \log \density(\btheta) + \textstyle \sum_{i = 1}^n \log \likelihood(y_i \given \btheta);
\end{tightEquation*}
where the sum of the $n$ i.i.d.\ terms should eventually dominate.

\pause
Here we will study some asymptotic behaviors of posteriors.
\end{frame}


\begin{frame}
\frametitle{Measuring ``distance'' between prob densities}
\textit{Kullback–Leibler divergence} measure how far away a probability density $p(\cdot)$ is from the reference $q(\cdot)$ and is defined as
\begin{tightEquation*}
\kldivergence( p \| q)
	\defeq - \int \log\!\left( \frac{q(x)}{p(x)} \right) p(x) \diff x
	= \expectation_{X \thinnestspace \sim \, p(\cdot)}\!\left[ \thinnerspace \log\!\left( \frac{q(X)}{p(X)} \right) \right].
\end{tightEquation*}
\pause%
\noteBullet To see relevance, think $p(\cdot) = \likelihood(\thinnerspace \cdot \given \btheta)$ and $q(\cdot) = \likelihood(\thinnerspace \cdot \given \btheta_\truthSub)$.

\pause
\smallskip
In particular, the KL divergence has the property 
\begin{tightEquation*}
\kldivergence(p \| q) \geq 0
	\ \text{ and } \
	\kldivergence(p \| q) = 0 \, \text{ only if } \, p(\cdot) \equiv q(\cdot).
\end{tightEquation*}
%$\kldivergence(p \| q) \geq 0$ and $\kldivergence(p \| q) = 0$ only if $p(\cdot) \equiv q(\cdot)$. 

\pause
\textit{Proof:} 
Jensen's inequality says $f(\expectation[X]) \leq \expectation[f(X)]$ for convex $f$, with the equality only if $f(x) = x$ on the support of $\density_X(\cdot)$. Hence,
\begin{tightEquation*}
0 = - \log\!\left( \int q(x) \diff x \right) \pause
	= - \log\!\left( \int \frac{q(x)}{p(x)} p(x) \diff x \right) \pause
	\leq \kldivergence(p \| q).
\end{tightEquation*}
\end{frame}


\begin{frame}
\frametitle{Posterior concentration around true parameter}

We now show that posteriors necessarily concentrate around the true parameter $\btheta_\truthSub$ as $n \to \infty$. 
To see this, we observe that
\smallskip
\begin{tightEquation*}
\begin{aligned}
\log\!\left( 
	\frac{\density(\btheta_\truthSub \given \by^n)}{\density(\btheta \given \by^n)}
\right)
	&= \log\!\left( 
		\frac{%
			\density(\btheta_\truthSub) \prod_i \likelihood(y_i \given \btheta_\truthSub)
		}{%
			\density(\btheta) \prod_i \likelihood(y_i \given \btheta)
		}%
	\right) \\
	\pause
	&= \log\!\left( 
		\frac{%
			\density(\btheta_\truthSub)
		}{%
			\density(\btheta)
		}%
	\right)
	+ \sum_{i = 1}^n \log\!\left( 
		\frac{%
			\likelihood(y_i \given \btheta_\truthSub)
		}{%
			\likelihood(y_i \given \btheta)
		}%
	\right).
\end{aligned}
\end{tightEquation*}
\pause%
\smallskip
Now we only need to recognize the latter term as a KL divergence:
\smallskip
\begin{tightEquation*}
\frac{1}{n} \sum_i \log\!\left( 
	\frac{%
		\likelihood(y_i \given \btheta_\truthSub)
	}{%
		\likelihood(y_i \given \btheta)
	}%
\right) 
	\to - \int \log\!\left( 
		\frac{%
			\likelihood(y \given \btheta_\truthSub)
		}{%
			\likelihood(y \given \btheta)
		}%
	\right) \likelihood(y \given \btheta_\truthSub) \diff y.
\end{tightEquation*}

\pause%
Thus
$\displaystyle 
	\frac{\density(\btheta_\truthSub \given \by^n)}{\density(\btheta \given \by^n)}
\to \infty$ whenever $\btheta \neq \btheta_\truthSub$;
i.e.\ $\density(\thinnerspace \cdot \given \by^n) \to \delta_{\btheta_\truthSub}(\cdot)$.
	% Note: as long as the parameter is identifiable.
\end{frame}

\begin{frame}
\frametitle{Posterior concentration under mis-specified model}
Analogous result holds when a model is \textit{mis-specified} ---
the family $\{ \likelihood(\thinnerspace \cdot \given \btheta) \}_{\btheta}$ doesn't include the true data dist $y \sim \likelihood_\truthSub(\cdot)$; e.g.\
\begin{narrowItemize}[<+->]
\item $\operatorname{Skewness}(y) \neq 0$, but $\likelihood(\thinnerspace \cdot \given \mu, \phi) \sim \normalDist(\mu, \phi^{-1})$.
\item $\expectation[y] < 0$, but assumed $\mu \geq 0$ under $\likelihood(\thinnerspace \cdot \given \mu, \phi) \sim \normalDist(\mu, \phi^{-1})$.
\item $y = f(\bx) + \sigma(\bx) \epsilon$ for non-linear $f$ and non-const $\sigma$, \\ but $\likelihood(\thinnerspace \cdot \given \bx, \bbeta, \phi) \sim \normalDist(\bx^\transpose \bbeta, \phi^{-1})$.
\end{narrowItemize}

\pause[\thebeamerpauses]
Under mis-specification, we can still speak of an optimal approx
\begin{tightEquation*}
\btheta_{\mathrm{opt}}
	= \textstyle \argmin_{\btheta} \kldivergence\!\left(
		\likelihood(\thinnerspace \cdot \given \btheta) \, \| \, \likelihood_\truthSub
	\right),
\end{tightEquation*}
\pause%
\mbox{and similarly show $\displaystyle \frac{\density(\btheta_{\mathrm{opt}} \given \by^n)}{\density(\btheta \given \by^n)}
\to \infty$ if $\btheta_{\mathrm{opt}}\!$ is the unique minimizer.}
\end{frame}


\begin{frame}
\frametitle{Asymptotic Gaussianity of posteriors}
Besides the posterior concentration, we in fact have a Bayesian asymptotic theory analogous to the frequentist one for {\small MLE}'s:

\pause
\begin{theorem}[Bernstein-von-Mises/Bayesian Central Limit]
Under an appropriate regularity assumption, \\
\smallskip%
\makebox[\linewidth]{
$\density(\thinnerspace \cdot \given \by^n) 
	\to \normalDist\!\left(\btheta_\truthSub, \left[
		 n \infoMat(\btheta_\truthSub)
\right]^{-1} \right)$ \ as $\, n \to \infty$
}\\ % Note: This can be more precisely stated as the convergence of $n^{-1/2} (\Theta - \btheta_\truthSub)$.
\smallskip%
where 
$\, \infoMat(\btheta) = \expectation_{\thinnerspace y \thinnerspace \sim \thinnerspace  \likelihood(\thinnerspace \cdot \given \btheta)\!}\!\left[ 
	- \frac{\partial^2}{\partial \btheta^2} \log \likelihood(y \given \btheta)
\right]$.
\end{theorem}

\pause
\textbf{Cf.} Asymptotic theory for {\small MLE}: 
$\, \mle{\btheta} \to \normalDist(\btheta_\truthSub, \left[
		 n \infoMat(\btheta_\truthSub)
\right]^{-1})$.

\pause
\mbox{In particular, Bayes and freq CI's necessarily coincide asymptotically}
(even though their interpretations are completely distinct).
\end{frame}


\newcommand{\samplesizeDependentMap}{\widehat{\theta}^{\text{\raisebox{.1ex}{$\thinnerspace n$}}}}
\begin{frame}
\frametitle{Asymptotic Gaussianity: Proof sketch (univariate)}
Consider the Taylor expansion of $\log \density(\theta \given \by^n)$ around its mode $\samplesizeDependentMap$:
\begin{tightEquation*}
\begin{aligned}
\log \density(\theta \given \by^n)
	= \log \density\big( \samplesizeDependentMap \given \by^n \big)
	&+ \frac12 \big( \theta - \samplesizeDependentMap \big)^2 \frac{\diff^2}{\diff \theta^2} \log \density\big( \samplesizeDependentMap \given \by^n \big) \\
	&\hspace*{-1.5em}+ \frac{1}{3!} \big( \theta - \samplesizeDependentMap \big)^3 \frac{\diff^3}{\diff \theta^3} \log \density\big( \samplesizeDependentMap \given \by^n \big) 
		+ \ldots.
\end{aligned}
\end{tightEquation*}
\pause%
Observe that 
\begin{tightEquation*}
\begin{aligned}
\frac1n \frac{\diff^k}{\diff \theta^k} \log \density\big( \theta \given \by^n \big) 
	&= \frac1n \frac{\diff^k}{\diff \theta^k} \log \density(\theta) 
		+ \frac1n \sum_{i = 1}^n \frac{\diff^k}{\diff \theta^k} \log \likelihood(y_i \given \theta) \\
	\pause%
	&\to \expectation_{\thinnerspace y \thinnerspace \sim \thinnerspace  \likelihood(\thinnerspace \cdot \given \theta_\truthSub)\!} \!\left[
		\frac{\diff^k}{\diff \theta^k} \log \likelihood(y \given \theta)
	\right] \, \text{ as } \thinnerspace n \to \infty.
\end{aligned}
\end{tightEquation*}
\pause%
Since the posterior concentrates around $\samplesizeDependentMap \approx \theta_\truthSub$, the terms $\big( \theta - \samplesizeDependentMap \big)^3 \frac{\diff^3}{\diff \theta^3} ...$ becomes immaterial relative to $n \big( \theta - \samplesizeDependentMap \big)^2$. 
\pause%
\hfill \qedsymbol
\end{frame}

\begin{frame}
\frametitle{``Counterexample'' to the asymptotic theory}
In invoking the law of large numbers in the proof, the fixed dimen- sionality of $\btheta$ and i.i.d.\ nature of $y_i$'s were important assumptions.

\pause%
B{\small UT} these assumptions are violated by many modern statistical \mbox{methods; e.g.\ random effects, model selection, non-parametrics, etc.}

\pause%
\mbox{(In fact, you should go Bayes only when classical methods fall short.)}
\end{frame}


\begin{frame}
\frametitle{``Counterexample'' to the asymptotic theory}
One counterexample comes from a simple random effects model
\begin{tightEquation*}
y_{i1}, y_{i2} \given \mu_i, \phi
	\sim \normalDist(\mu_i, \phi^{-1}), \ \,
\mu_i \sim \normalDist(\mu_{\textrm{grp}}, \phi_{\textrm{grp}}^{-1})
\end{tightEquation*}
for $i = 1, \ldots, n$ with ``across-group'' mean $\mu_{\textrm{grp}}$ and variance $\phi_{\textrm{grp}}^{-1}$.
% Note: For simplicity, we take the across-group mean and variance to be fixed.

\pause%
(This is basically the model behind the Neyman-Scott paradox except for slight extension to incorporate shrinkage toward $\mu_{\textrm{grp}}$.)

\pause%
\smallskip
For example, $y_{i1}$ and $y_{i2}$ could represent batting averages of the $i$-th player in his 1st and 2nd games (\'{a} la \citealt{efron1975stein_est}).

\pause%
\smallskip
\textbf{Question:} Would the posterior of $(\bmu, \phi)$ concentrate around the truth for $n$ large? 
Would the posterior look Gaussian?
\end{frame}


\begin{frame}
\frametitle{``Counterexample'' to the asymptotic theory}
For simplicity, let's just assume the objective prior $\density(\bmu, \phi) \propto 1 / \phi$ \\
(i.e.\ $\phi_{\textrm{grp}} = \infty$), ignoring the shrinkage estimation aspect.

\pause%
Asymptotic behavior, if real, shouldn't depend on a prior anyways.

\pause%
The corresponding posterior is
\begin{tightEquation*}
\begin{gathered}
\mu_i \given \bm{y}, \phi 
	\sim \normalDist\!\left( 
		\frac{y_{i1} + y_{i2}}{2}, \thinnerspace
		\frac12 \phi^{-1}
	\right), \\
\phi \given \bm{y}
	\sim \gammaDist\!\left( 
		\frac{n}{2}, \thinnerspace 
		\frac{\sum_i (y_{i1} - y_{i2})^2}{4} 
	\right).
\end{gathered}
\end{tightEquation*}

\pause%
The posterior means provide unbiased estimates of the truths:
\begin{tightEquation*}
\expectation[\mu_i \given \bm{y}]
%	= \expectation[ \expectation[\mu_i \given \bm{y}, \phi] ]
	= \frac{y_{i1} + y_{i2}}{2}, \ \,
\expectation[\phi \given \bm{y}]^{-1}
	= \frac{\sum_i (y_{i1} - y_{i2})^2}{2 n}.
\end{tightEquation*}

\pause%
In fact, $\density(\phi \given \by) \to \delta_{\phi_\truthSub}(\phi)$, though $\variance(\mu_i \given \bm{y})$ never vanishes.
\end{frame}


\begin{frame}
\frametitle{``Counterexample'' to the asymptotic theory}
We already see the asymtotic theory breaking down, but let's nonetheless ask the question: Does the posterior look Gaussian?

\pause%
\smallskip
That in particular would imply its mean and mode to be similar.

\pause%
\smallskip
But the posterior mode exhibits a fundamentally different behavior:
%But the posterior mode (which is similar to the {\small MLE}) exhibits a fundamentally different behavior, being off by a factor of 2:
\begin{tightEquation*}
\map{\phi}^{-1} = \frac{\sum_i (y_{i1} - y_{i2})^2}{4 (n - 1)}.
\end{tightEquation*}
% Note: Worth working this out in the class for students to really believe it.

\pause%
The mode, and {\small MLE}, is off by a factor of 2! (Hence a ``paradox.'')
\end{frame}


\begin{frame}
\frametitle{``Counterexample'' to the asymptotic theory}
Why the difference? 

\pause%
The answer, turns out, is not so much ``mean vs.\ mode'' but the Bayes estimator properly accounting for uncertainty in $\mu_i$'s. 

\pause%
\smallskip
To be more precise, the mean is based on the posterior marginal
\begin{tightEquation*}
\density(\phi \given \by)
	= \int \density(\phi \given \bmu, \by) \thinnerspace \density(\bmu \given \by) \diff \bmu
\end{tightEquation*}
which \textit{integrates out} $\bmu$, whose uncertainty remains as $n \to \infty$.

\pause%
\smallskip
%Incidentally, the posterior marginal can be equivalently derived as $\propto \likelihood(\by \given \phi) \density(\phi)$ in terms of the marginal likelihood
%\begin{tightEquation*}
%\likelihood(\by \given \phi)
%%	= \int \likelihood(\by, \bmu \given \phi) \diff \bmu
%	= \int \likelihood(\by \given \bmu, \phi) \density(\bmu \given \phi) \diff \bmu.
%\end{tightEquation*}
Its effect is also seen in the fact $\thinnerspace \variance(\phi \given \by) \sim 4 n / \phi_\truthSub^2$ while 
\begin{tightEquation*}
- \frac{\diff^2}{\diff (\bmu, \phi)^2} \log \likelihood(y \given \bmu, \phi)
	= \begin{bmatrix}
	\phi \Id_{n \times n} & \bm{0} \\
	0 & n / \phi^2
	\end{bmatrix}.
\end{tightEquation*}
\end{frame}


\section{\centerline{Summary}}

\begin{frame}
\frametitle{Summary of ``Bayes under parametric models''}
We have studied behaviors of posteriors under relatively simple models to build intuitions. 

\pause%
Under parametric models, Bayesian inference has much similarity to the classical one based on {\small MLE}.
\pause%
\begin{narrowItemize}[<+->]
\item Conjugate analysis demonstrates sound, intuitive behaviors of posteriors and helps quantify relative weights on priors. 
	% Note: Having been trained in the classical statistics, we cannot help notice the familiar quantities popping up in the Bayesian analysis of the models. But the only thing we did really was just to adhere to what the probability theory tells us and multiply the likelihood and prior. We woud have arrived at a right statistical inference even if we knew nothing about the classical statistical estimation methods.
	% Note: Many Bayes newbies are, understandably, uncomfortable at picking a prior. But the prior really doesn't have much impact on the posterior when the likelihood is informative and, when the likelihood is not, then we cannot really proceed without imposing some additional structures.
\item Bayesian cred intervals often have remarkably good freq performance (despite the Bayes-freq philosophical differences).
\item The two``CI''s coincide asymptotically; 
	Bayes's may well have better freq performance than alternatives in finite samples.
%\item In fact, Bayesian thinking can be useful in devising a freq procedure, whose performance can then be evaluated formally.
\end{narrowItemize}
\end{frame}

\nobibliography{references}
\newcommand{\bda}{BDA3}
\begin{frame}
\frametitle{References}

Simple parametric models:

\begin{indented}[4]
\bda{} Chapter~2 \& 3 for uni-/multi-variate models.
\end{indented}
%\begin{narrowItemize}
%\item \bda{} Chapter~2/3 for uni-/multi-variate models
%\end{narrowItemize}

Linear regression:
\begin{narrowItemize}
\item \bda{} Chapter~14.1--14.5.
\end{narrowItemize}

Bayesian asymptotics:
\begin{narrowItemize}
\item \bda{} Chapter~4 and Appendix~B.
\end{narrowItemize}

Further readings:
\begin{narrowItemize}
\item \cite{bayarri2004bayes_freq_interplay}.% 
	``\usebibentry{bayarri2004bayes_freq_interplay}{title}.''%
	\textit{\usebibentry{bayarri2004bayes_freq_interplay}{journal}}.
\end{narrowItemize}
\end{frame}


\end{document}
